%% Apéndice: notación. Recapitulación; los signos se introducen en la Parte I.
\chapter{Notación}
\label{ap:notacion}

\section{Los signos}


\begin{center}\small
\begin{tabularx}{\linewidth}{@{}l >{\raggedright\arraybackslash}X@{}}
\toprule
signo & qué es \\
\midrule
$\defnot{\Gamma \vdash A}$ & $A$ se deriva de las hipótesis $\Gamma$ en el cálculo $\omega$ \\
$\defnot{\Prf\,A}$ & $A$ se deriva en el cálculo de Hilbert finitario, sin hipótesis \\
$\sigma n$ & el sucesor de $n$ \\
$t \doteq u$ & igualdad \textbf{del lenguaje objeto} (no la de Lean, que es \ident{=}) \\
\texttt{\#0} & la variable ligada más cercana: un \defterm{índice de De Bruijn} \\
$\defnot{\codigo{\varphi}}$ & el \defterm{código de Gödel} de $\varphi$: un término objeto que la denota \\
$\defnot{\dotado{a}}$ & «$a$ dotado»: el código de $a$ construido \emph{dentro} del lenguaje (\ident{tcFn}) \\
$\defnot{\Prov(\codigo{\varphi})}$ & la fórmula objeto que dice «$\varphi$ es demostrable» \\
$\Delta_0$, $\Sigma_1$ & clases de fórmulas por su cuantificación: acotada, y $\exists$ sobre acotada \\
\bottomrule
\end{tabularx}
\end{center}

\nota{Ningún término técnico aparece en este libro antes de estar definido aquí. Lo comprueba
\modulo{doc/book/scripts/terminos.py} en cada compilación, y no admite excepciones.}

\section{Cómo se lee un nombre}

Los identificadores van en inglés y siguen las convenciones de Mathlib. Tres reglas bastan para
leer casi cualquiera:

\begin{enumerate}[leftmargin=1.4em]
\item \textbf{La conclusión primero, las hipótesis con \ident{_of_}.} El nombre dice \emph{qué se
  demuestra}, no cómo. \ident{isNat_succ_of_isNat} es «$\sigma n$ es natural, \emph{a partir de}
  que $n$ lo es».
\item \textbf{Los bicondicionales acaban en \ident{_iff}.} \ident{prf_In_iff_boundedIn} es un «si y
  sólo si».
\item \textbf{Los símbolos se deletrean}: \ident{add}, \ident{mul}, \ident{succ}, \ident{eq},
  \ident{lt}, \ident{le}, \ident{dvd}, \ident{not}, \ident{forall}, \ident{exists}. Así
  \ident{prf_add_zero_left} se lee sin traducir.
\end{enumerate}

Y hay algo más útil todavía, porque conecta con la chapa de niveles:
\textbf{el prefijo dice en qué nivel se afirma.} No es una impresión — es la práctica medida del
proyecto:

\begin{center}\small
\begin{tabularx}{\linewidth}{@{}l r >{\raggedright\arraybackslash}X@{}}
\toprule
prefijo & símbolos & qué anuncia \\
\midrule
\ident{ax_}   & 154 & un \defterm{axioma objeto}: se postula, no se demuestra \\
\ident{prf_}  & 585 & un teorema del cálculo finitario \ident{Prf} \\
\ident{pcc_}  & 214 & un teorema de \ident{Prf} \textbf{dentro de} \ident{Prov} — la teoría hablando de sí misma \\
\ident{CRIT_} & 13  & un \defterm{sondeo} crítico: una pregunta contestada, a menudo en negativo \\
\bottomrule
\end{tabularx}
\end{center}

\nota{Los recuentos son de \modulo{ROBINSON\_PlusPlus/} a 2026-09-03. La excepción conocida son las
seis meta-reglas $\omega$ de \modulo{FOL}, anteriores a la convención, que no llevan \ident{ax_}.}

\nota{Ningún término ni notación aparece en este libro antes de haberse introducido. Lo comprueba
\modulo{doc/book/scripts/terminos.py} en cada compilación, y no admite excepciones.}
